\chapter{HỌC TĂNG CƯỜNG SÂU VỚI BIỂU DIỄN HÀNH ĐỘNG CÓ CẤU TRÚC VẬT LÝ}
\label{chap:phuong-phap}

\section{Tổng quan khung phương pháp}
\label{sec:tong-quan-khung}

Mục tiêu của Chương~\ref{chap:phuong-phap} là chuyển bài toán tối ưu ở
Chương~\ref{chap:mo-hinh-he-thong} thành một quy trình có thể học. Khung phương pháp gồm
bốn lớp. Lớp thứ nhất nhận CSI và chuyển thành véc-tơ trạng thái có chuẩn hóa. Lớp thứ
hai là thuật toán DRL, tạo véc-tơ hành động chuẩn hóa trong miền $[-1,1]$. Lớp thứ ba là
bộ giải mã hành động vật lý, chuyển véc-tơ chuẩn hóa thành các biến $\mathbf{p}$,
$\boldsymbol{\eta}$, $\boldsymbol{\beta}$, $\boldsymbol{\theta}^T$,
$\boldsymbol{\theta}^R$ thỏa các ràng buộc cơ bản. Lớp cuối là môi trường
STAR--RIS--RSMA, tính kênh hiệu dụng, tốc độ, QoS và phần thưởng.

Điểm quan trọng là bộ giải mã dùng chung cho DDPG, TD3 và PPO. Vì vậy, khi so sánh ba
thuật toán, sự khác nhau chủ yếu nằm ở cách chúng cập nhật chính sách chứ không nằm ở
việc một thuật toán được cấp một miền hành động thuận lợi hơn thuật toán khác. Đây là
nguyên tắc cần thiết để đánh giá công bằng vai trò của nhóm học ngoài chính sách và nhóm
học theo chính sách hiện hành.

\begin{figure}[htbp]
\centering
\resizebox{0.95\textwidth}{!}{\begin{tikzpicture}[font=\scriptsize,>=Stealth,line width=0.5pt]

\node[bx,text width=2.0cm,minimum height=1.0cm] (csi)  at (0,0)    {CSI\\ $h_d,\ g,\ h_r$};
\node[bx,text width=2.2cm,minimum height=1.0cm] (st)   at (3.0,0)  {Tạo trạng thái\\ và chuẩn hóa\\ \tiny $s\in\mathbb{R}^{d_s}$};
\node[bx,text width=2.6cm,minimum height=1.0cm] (net)  at (6.4,0)  {Actor\\ DDPG / TD3 / PPO};
\node[bx,text width=2.4cm,minimum height=1.0cm] (act)  at (9.8,0)  {Hành động chuẩn hóa\\ $a\in[-1,1]^{d_a}$};
\node[bx,fill=cChung!10,text width=2.6cm,minimum height=1.0cm] (dec) at (9.8,-2.3)
      {Bộ giải mã hành động vật lý\\ \tiny $p,\ \eta,\ \beta,\ \theta^T,\ \theta^R$};
\node[bx,fill=cKenh!10,text width=2.8cm,minimum height=1.0cm] (env) at (5.4,-2.3)
      {Môi trường\\ STAR--RIS--RSMA};
\node[bx,text width=2.4cm,minimum height=1.0cm] (rw)  at (1.4,-2.3)
      {Tốc độ, QoS\\ và phần thưởng};

\draw[->] (csi)  -- (st);
\draw[->] (st)   -- (net);
\draw[->] (net)  -- (act);
\draw[->] (act)  -- (dec);
\draw[->] (dec)  -- (env);
\draw[->] (env)  -- (rw);
\draw[->] (rw.north) -- (1.4,-1.15) -- (6.4,-1.15) -- (net.south);
\node[font=\tiny,anchor=south] at (3.9,-1.12) {tín hiệu học (cập nhật trọng số)};

\node[note,text width=12.5cm] at (-0.9,-3.30)
  {Bộ giải mã hành động vật lý \textbf{dùng chung} cho DDPG, TD3 và PPO, nên khác biệt
   giữa ba thuật toán chỉ nằm ở cách cập nhật chính sách, không nằm ở miền hành động.};
\end{tikzpicture}
}
\caption{Khung DRL và bộ giải mã hành động vật lý dùng chung}
\label{fig:ch3-khung-drl}
\figuresource{Tác giả xây dựng theo mã nguồn thực nghiệm (2026)}
\end{figure}

\section{Biểu diễn trạng thái}
\label{sec:bieu-dien-trang-thai}

\subsection{Cấu trúc trạng thái}
\label{ssec:cau-truc-trang-thai}

Trạng thái đưa vào mạng phải chứa đủ thông tin để tái tạo tác động của một hành động lên
tốc độ. Mã nguồn sử dụng phần thực và phần ảo của kênh trực tiếp, kênh BS--STAR--RIS,
kênh STAR--RIS--UE cùng chỉ báo phía của UE. Số chiều trạng thái là:
\begin{equation}
\label{eq:so-chieu-trang-thai}
d_s = 2(K + N + KN) + K.
\end{equation}

Trong đó: $K$ là số UE; $N$ là số phần tử STAR--RIS; hệ số $2$ xuất hiện vì mỗi hệ số
phức được tách thành phần thực và phần ảo; $K$ cuối cùng là chỉ báo phía của các UE. Phần
lớn thông tin trạng thái đến từ ma trận kênh STAR--RIS--UE có $KN$ phần tử. Khi $N$ tăng,
mạng nhận nhiều thông tin hơn nhưng vẫn giữ cấu trúc véc-tơ cố định cho từng cấu hình.
Với $K = 4$ và $N = 32$, số chiều trạng thái bằng $332$.

\subsection{Chuẩn hóa theo khối}
\label{ssec:chuan-hoa-theo-khoi}

Ba nhóm kênh có mức tỷ lệ khác nhau ngay từ bước sinh dữ liệu. Vì vậy, trước khi đưa vào
mạng, từng nhóm được chia cho đúng hệ số đã dùng để sinh kênh:
\begin{equation}
\label{eq:chuan-hoa-khoi}
\tilde{h}_{d,k} = \frac{h_{d,k}}{0,35},
\qquad
\tilde{g}_n = \frac{g_n}{1,0},
\qquad
\tilde{h}_{r,k,n} = \frac{h_{r,k,n}}{0,75}.
\end{equation}

Trong đó: dấu ngã chỉ đại lượng sau chuẩn hóa; các mẫu số $0,35$, $1,0$ và $0,75$ là đúng
các hệ số tỷ lệ trong hàm sinh kênh. Phép chia không làm thay đổi thứ tự mạnh yếu tương
đối bên trong mỗi khối, mà đưa các khối về cùng bậc độ lớn. Nếu không chuẩn hóa, khối có
biên độ lớn hơn có thể chi phối gradient của lớp đầu chỉ vì đơn vị số học chứ không phải
vì thông tin vật lý quan trọng hơn.

Sau chuẩn hóa, các giá trị được cắt trong khoảng $[-8, 8]$. Mục đích là hạn chế các mẫu
Gaussian hiếm có biên độ quá lớn làm mất ổn định quá trình học. Chỉ báo phía UE được ánh
xạ từ $\{0, 1\}$ sang $\{-1, 1\}$ để có thang gần với các thành phần còn lại.
Hình~\ref{fig:ch3-trang-thai} tóm tắt cấu trúc này: bốn khối của véc-tơ trạng thái với số
chiều ký hiệu, giá trị cụ thể khi $K = 4$ và $N = 32$, hệ số chuẩn hóa của từng khối và
bước cắt biên trước khi đưa vào mạng.

\begin{figure}[htbp]
\centering
\resizebox{\textwidth}{!}{\begin{tikzpicture}[font=\footnotesize,>=Stealth]
\def\h{0.85}
% cac khoi
\node[sbox,fill=blue!12,  minimum width=2.1cm,minimum height=\h cm,anchor=west] (b1) at (0,0)    {$\mathrm{Re}/\mathrm{Im}\ h_{d,k}$};
\node[sbox,fill=green!14, minimum width=2.6cm,minimum height=\h cm,anchor=west] (b2) at (2.1,0)  {$\mathrm{Re}/\mathrm{Im}\ g_n$};
\node[sbox,fill=orange!16,minimum width=5.2cm,minimum height=\h cm,anchor=west] (b3) at (4.7,0)  {$\mathrm{Re}/\mathrm{Im}\ h_{r,k,n}$};
\node[sbox,fill=violet!14,minimum width=2.3cm,minimum height=\h cm,anchor=west] (b4) at (9.9,0)  {chỉ báo phía $s_k$};

% so chieu duoi tung khoi
\node[font=\scriptsize,anchor=north] at (1.05,-0.50) {$2K$};
\node[font=\scriptsize,anchor=north] at (3.40,-0.50) {$2N$};
\node[font=\scriptsize,anchor=north] at (7.30,-0.50) {$2KN$};
\node[font=\scriptsize,anchor=north] at (11.05,-0.50) {$K$};
\node[font=\tiny,anchor=north,gray!60!black] at (1.05,-0.88) {$=8$};
\node[font=\tiny,anchor=north,gray!60!black] at (3.40,-0.88) {$=64$};
\node[font=\tiny,anchor=north,gray!60!black] at (7.30,-0.88) {$=256$};
\node[font=\tiny,anchor=north,gray!60!black] at (11.05,-0.88) {$=4$};

% he so chuan hoa phia tren
\node[font=\scriptsize,anchor=south] at (1.05,0.52) {$\div\,0{,}35$};
\node[font=\scriptsize,anchor=south] at (3.40,0.52) {$\div\,1{,}0$};
\node[font=\scriptsize,anchor=south] at (7.30,0.52) {$\div\,0{,}75$};
\node[font=\scriptsize,anchor=south] at (11.05,0.52) {$\{0,1\}\!\to\!\{-1,1\}$};

% ngoac tong
\draw[decorate,decoration={brace,amplitude=5pt,mirror},gray!70]
      (0,-1.35) -- (12.2,-1.35);
\node[font=\scriptsize,anchor=north] at (6.1,-1.80)
     {$d_s = 2(K+N+KN)+K = 332$ \quad với $K=4,\ N=32$};

% buoc cat bien
\node[box,anchor=west,font=\scriptsize] (clip) at (0,-2.75) {Cắt biên về $[-8,\,8]$};
\node[box,anchor=west,font=\scriptsize] (net)  at (4.2,-2.75) {Mạng nơ-ron (Actor / Critic)};
\draw[->] ($(clip.north)+(0,0.60)$) -- (clip.north);
\draw[->] (clip.east) -- (net.west);
\end{tikzpicture}
}
\caption{Cấu trúc và chuẩn hóa theo khối của véc-tơ trạng thái}
\label{fig:ch3-trang-thai}
\figuresource{Tác giả xây dựng theo công thức \eqref{eq:so-chieu-trang-thai}, \eqref{eq:chuan-hoa-khoi} và mã nguồn (2026)}
\end{figure}

\section{Không gian hành động chuẩn hóa}
\label{sec:khong-gian-hanh-dong}

Mạng không sinh trực tiếp công suất và pha. Đầu ra được tổ chức thành véc-tơ chuẩn hóa:
\begin{equation}
\label{eq:vec-to-hanh-dong}
a = \big[\, a_c,\ a_p,\ a_\eta,\ a_\beta,\ a_{\theta^T},\ a_{\theta^R} \,\big]
\in [-1,1]^{d_a}.
\end{equation}

Trong đó: $a_c$ điều khiển tỷ lệ công suất luồng chung; $a_p \in \mathbb{R}^{K}$ điều
khiển phân bổ công suất riêng; $a_\eta \in \mathbb{R}^{K}$ điều khiển phân bổ tốc độ
chung; ba khối còn lại có kích thước $N$ và điều khiển năng lượng, pha truyền, pha phản
xạ.

Cách tổ chức này tách ngôn ngữ của mạng nơ-ron khỏi ngôn ngữ vật lý. Mạng chỉ cần học các
số có cùng miền $[-1,1]$; bộ giải mã chịu trách nhiệm chuyển từng khối sang đúng ràng
buộc vật lý, nhờ đó giảm việc mạng phải tự học các phép chiếu và giới hạn. Số chiều của
$a$ chính là $d_a = 2K + 1 + 3N$ đã nêu ở Chương~\ref{chap:mo-hinh-he-thong}. Dù cách
giải mã có cấu trúc, số lượng đầu ra không giảm; thay đổi chính nằm ở ý nghĩa của các đầu
ra và hình học miền vật lý mà chúng có thể biểu diễn.

\section{Bộ giải mã hành động có cấu trúc vật lý}
\label{sec:bo-giai-ma}

\subsection{Giải mã công suất luồng chung}
\label{ssec:giai-ma-cong-suat-chung}

Đầu ra $a_c \in [-1,1]$ được ánh xạ tuyến tính thành tỷ lệ công suất luồng chung:
\begin{equation}
\label{eq:rho-c}
\rho_c = 0,70 + (0,99 - 0,70)\,\frac{a_c + 1}{2}.
\end{equation}

Trong đó: $\rho_c$ là tỷ lệ công suất dành cho luồng chung; khi $a_c = -1$ thì $\rho_c =
0,70$ và khi $a_c = 1$ thì $\rho_c = 0,99$. Ánh xạ này chủ động loại khỏi miền hành động
các tỷ lệ công suất luồng chung quá thấp. Phân tích Chương~\ref{chap:mo-hinh-he-thong}
cho thấy luồng chung là nguồn tốc độ quan trọng cho UE không được ưu tiên công suất luồng
riêng, vì vậy vùng công suất chung cao tạo một điểm xuất phát thuận lợi hơn cho QoS.
Khoảng này là lựa chọn thiết kế, không phải ràng buộc bắt buộc của RSMA.

Công suất luồng chung được tính:
\begin{equation}
\label{eq:p-c}
p_c = \rho_c P_{\max}.
\end{equation}

Trong đó: $P_{\max}$ là tổng ngân sách công suất và $p_c$ là phần được dành cho luồng
chung. Cách tính này bảo đảm công suất chung luôn không âm và không vượt ngân sách; phần
còn lại $(1-\rho_c)P_{\max}$ được chuyển toàn bộ sang nhóm luồng riêng, vì vậy tổng công
suất luôn được sử dụng hết.

\subsection{Giải mã công suất luồng riêng}
\label{ssec:giai-ma-cong-suat-rieng}

Trọng số luồng riêng được tính bằng softmax với hệ số độ sắc $\kappa = 10$:
\begin{equation}
\label{eq:softmax}
w_k = \frac{\exp(\kappa a_{p,k})}{\sum_{j=1}^{K} \exp(\kappa a_{p,j})},
\qquad \kappa = 10.
\end{equation}

Trong đó: $a_{p,k}$ là đầu ra chuẩn hóa cho UE $k$; $w_k > 0$ và $\sum_k w_k = 1$;
$\kappa$ là hệ số độ sắc, tương đương hệ số nghịch nhiệt độ nếu dùng cách gọi theo softmax
chuẩn. Softmax ở đây khuếch đại chênh lệch giữa các đầu ra Actor: chỉ cần một UE có
$a_{p,k}$ lớn hơn rõ rệt, $w_k$ có thể tiến rất gần $1$ trong khi các trọng số khác tiến
rất gần $0$. Nhờ đó mạng biểu diễn được phân bổ công suất luồng riêng gần như dồn hết vào
một luồng bằng các đầu ra hữu hạn, phù hợp trực giác SISO ở
Chương~\ref{chap:mo-hinh-he-thong}.

Công suất riêng được tính:
\begin{equation}
\label{eq:p-k}
p_k = (1 - \rho_c) P_{\max} w_k.
\end{equation}

Trong đó: $1 - \rho_c$ là tỷ lệ công suất còn lại sau khi cấp cho luồng chung và $w_k$
quyết định UE nào nhận bao nhiêu phần của ngân sách riêng. Cách tính này bảo đảm đồng
thời $p_k \ge 0$ và $\sum_k p_k = (1-\rho_c)P_{\max}$, nên kết hợp~\eqref{eq:p-c}
và~\eqref{eq:p-k} luôn thỏa đúng tổng công suất mà không cần thêm phép phạt hoặc chuẩn
hóa sau cùng.

Ví dụ, với $K = 4$ và $a_p = [1, -1, -1, -1]$, chênh lệch logit sau khi nhân $\kappa$ bằng
$20$. Tỷ lệ giữa trọng số lớn nhất và mỗi trọng số nhỏ vào khoảng $e^{20}$, nên phân bổ có
thể cực kỳ tập trung dù softmax về mặt toán học vẫn không tạo giá trị $0$ tuyệt đối. Do đó
luận văn dùng thuật ngữ gần one-hot thay vì one-hot chính xác.

Hình~\ref{fig:ch3-do-sac} cho thấy ảnh hưởng của hệ số $\kappa$ khi giữ nguyên một đầu ra
Actor. Với $\kappa = 1$, bốn trọng số gần như chia đều. Với $\kappa = 5$, UE có đầu ra lớn
nhất đã chiếm khoảng ba phần tư ngân sách riêng. Với $\kappa = 10$ như cấu hình của luận
văn, tỷ lệ đó vượt $0,94$ trong khi UE yếu nhất gần như bằng không. Đây là lý do thực tế
cho việc chọn hệ số độ sắc cao.

\begin{figure}[htbp]
\centering
\resizebox{0.95\textwidth}{!}{\begin{tikzpicture}[font=\footnotesize]
\begin{axis}[
  width=13.5cm, height=6.4cm,
  ybar=2pt, bar width=11pt,
  xlabel={Chỉ số UE $k$}, ylabel={Trọng số công suất riêng $w_k$},
  symbolic x coords={UE1,UE2,UE3,UE4},
  xtick=data, ymin=0, ymax=1.08,
  grid=both, grid style={gray!25},
  tick label style={font=\scriptsize}, label style={font=\scriptsize},
  legend style={font=\scriptsize,at={(0.98,0.97)},anchor=north east,draw=gray!50},
  legend cell align=left,
  nodes near coords, every node near coord/.append style=
     {font=\tiny,/pgf/number format/.cd,fixed,precision=3,use comma},
]
\addplot[fill=gray!35,draw=gray!60!black]      coordinates {(UE1,0.358)(UE2,0.265)(UE3,0.217)(UE4,0.161)};
\addlegendentry{$\kappa=1$}
\addplot[fill=orange!45,draw=orange!70!black]  coordinates {(UE1,0.756)(UE2,0.169)(UE3,0.062)(UE4,0.014)};
\addlegendentry{$\kappa=5$}
\addplot[fill=blue!45,draw=blue!70!black]      coordinates {(UE1,0.946)(UE2,0.047)(UE3,0.006)(UE4,0.000)};
\addlegendentry{$\kappa=10$ (luận văn)}
\end{axis}
\node[font=\scriptsize,anchor=north] at (6.8,-1.35)
     {Cùng một đầu ra Actor $a_p=[0{,}5;\ 0{,}2;\ 0{,}0;\ -0{,}3]$};
\end{tikzpicture}
}
\caption{Ảnh hưởng của hệ số độ sắc $\kappa$ đến phân bổ công suất luồng riêng}
\label{fig:ch3-do-sac}
\figuresource{Tác giả xây dựng theo công thức \eqref{eq:softmax} (2026)}
\end{figure}

\subsection{Giải mã tỷ lệ phân bổ tốc độ chung}
\label{ssec:giai-ma-ty-le-toc-do}

Khối $a_\eta$ được đưa từ $[-1,1]^K$ về $[0,1]^K$ rồi chiếu Euclid lên đơn hình:
\begin{equation}
\label{eq:eta}
\boldsymbol{\eta} = \Pi_{\Delta_K}\!\left(\frac{a_\eta + 1}{2}\right),
\qquad
\Delta_K = \Bigl\{\, x \;\big|\; x_k \ge 0,\ \textstyle\sum_{k} x_k = 1 \,\Bigr\}.
\end{equation}

Trong đó: $\Pi_{\Delta_K}$ là phép chiếu lên đơn hình xác suất $\Delta_K$ và
$\boldsymbol{\eta}$ là véc-tơ tỷ lệ phân bổ tốc độ chung. Phép chiếu biến một đầu ra tùy ý
của mạng thành một quy tắc chia toàn bộ quỹ tốc độ chung: nếu Actor muốn hỗ trợ UE yếu, nó
tăng thành phần $a_{\eta,k}$ tương ứng, và phép chiếu tự điều chỉnh các thành phần còn lại
để tổng vẫn bằng một.

\subsection{Giải mã hệ số chia năng lượng}
\label{ssec:giai-ma-chia-nang-luong}

Mỗi phần tử STAR--RIS dùng ánh xạ tuyến tính:
\begin{equation}
\label{eq:beta}
\beta^T_n = \frac{a_{\beta,n} + 1}{2},
\qquad
\beta^R_n = 1 - \beta^T_n.
\end{equation}

Trong đó: $a_{\beta,n} \in [-1,1]$ là đầu ra Actor; $\beta^T_n$ và $\beta^R_n \in [0,1]$
là tỷ lệ năng lượng truyền và phản xạ của phần tử thứ $n$. Ràng buộc bảo toàn năng lượng
được mã hóa trực tiếp: Actor chỉ cần quyết định nghiêng phần tử về phía truyền hay phía
phản xạ, còn bộ giải mã tự bảo đảm hai tỷ lệ cộng bằng một. Khi $a_{\beta,n} = 0$, năng
lượng được chia đều $0,5/0,5$.

\subsection{Pha tham chiếu phụ thuộc kênh}
\label{ssec:pha-tham-chieu}

Trọng số luồng riêng $w_k$ không chỉ dùng cho công suất mà còn chỉ ra UE nào đang được
chính sách ưu tiên. Bộ giải mã dùng các trọng số này để tạo tổng phức trên từng phần tử:
\begin{equation}
\label{eq:z-n}
z_n = \sum_{k=1}^{K} w_k\, h^{*}_{r,k,n}\, g_n .
\end{equation}

Trong đó: $h^{*}_{r,k,n}$ là liên hợp của kênh STAR--RIS--UE; $g_n$ là kênh
BS--STAR--RIS; $w_k$ là mức ưu tiên luồng riêng của UE. Tổng phức này tạo một hướng pha
đại diện của các UE được Actor ưu tiên: nếu một $w_k$ gần một, tổng gần như chỉ còn kênh
ghép tầng của UE đó; nếu nhiều trọng số đáng kể, pha tham chiếu là sự kết hợp có trọng số.

Pha tham chiếu là:
\begin{equation}
\label{eq:theta-ref}
\bar{\theta}_n = -\angle z_n .
\end{equation}

Trong đó: $\angle z_n$ là góc pha của số phức $z_n$ và dấu âm dùng để bù pha của đường
ghép tầng. Pha tham chiếu đóng vai trò tương tự một bước căn pha giải tích: trước khi mạng
học, mỗi phần tử đã được đặt quanh một hướng có khả năng làm các thành phần ghép tầng cộng
thuận cho hỗn hợp UE được ưu tiên, nhờ đó giảm gánh nặng phải khám phá toàn bộ vòng pha
$2\pi$ từ đầu.

\subsection{Học phần hiệu chỉnh pha}
\label{ssec:hieu-chinh-pha}

Actor không bị khóa cứng vào pha tham chiếu mà học hai phần hiệu chỉnh độc lập:
\begin{align}
\theta^T_n &= \mathrm{wrap}\!\left(\bar{\theta}_n + 0,25\pi\, a_{\theta^T\!,n}\right),
\label{eq:theta-T}\\[2pt]
\theta^R_n &= \mathrm{wrap}\!\left(\bar{\theta}_n + 0,25\pi\, a_{\theta^R\!,n}\right).
\label{eq:theta-R}
\end{align}

Trong đó: $\mathrm{wrap}(\cdot)$ đưa pha về miền $[-\pi, \pi)$; $a_{\theta^T\!,n}$ và
$a_{\theta^R\!,n}$ là hai đầu ra độc lập; biên $0,25\pi$ tương ứng phần hiệu chỉnh tối đa
$\pm\pi/4$.

Hai công thức kết hợp kiến thức giải tích và khả năng học. Pha tham chiếu xử lý phần căn
chỉnh thô, còn mạng học phần sửa sai để tính đến đường trực tiếp, nút cổ chai của luồng
chung, QoS, phân chia năng lượng và tương tác giữa hai phía. Hai phần hiệu chỉnh độc lập
giúp phía truyền và phía phản xạ có thể lệch khác nhau quanh cùng một tham chiếu.

Hình~\ref{fig:ch3-bo-giai-ma} tổng hợp năm phép giải mã vừa trình bày. Mỗi hàng của hình
đi từ một đoạn của véc-tơ hành động, qua phép biến đổi tương ứng, đến biến vật lý cùng
ràng buộc mà nó luôn thỏa. Hai khối pha dùng chung một bộ giải mã vì cùng lấy
$\bar{\theta}_n$ làm điểm neo.

\begin{figure}[htbp]
\centering
\resizebox{\textwidth}{!}{\begin{tikzpicture}[font=\footnotesize,>=Stealth]

% ---------- vec-to hanh dong, xep doc ----------
\node[font=\scriptsize,anchor=south] at (0.6,0.45) {$a\in[-1,1]^{d_a}$};
\foreach \y/\lbl/\dim in {0/{$a_c$}/1, -1.1/{$a_p$}/K, -2.2/{$a_\eta$}/K,
                          -3.3/{$a_\beta$}/N, -4.40/{$a_{\theta^T}$}/N, -5.05/{$a_{\theta^R}$}/N}{
   \node[sbox,fill=gray!18,minimum width=1.2cm,minimum height=0.55cm,anchor=west] at (0,\y) {\lbl};
   \node[font=\tiny,gray!55!black,anchor=east] at (-0.15,\y) {\dim};
}
\node[font=\scriptsize,anchor=north] at (0.6,-5.55) {$d_a=2K+1+3N$};

% ---------- khoi giai ma ----------
\node[box,anchor=west,text width=4.9cm,font=\scriptsize,minimum height=0.8cm] (d1) at (2.4,0)
     {Ánh xạ tuyến tính\\ $\rho_c=0{,}70+0{,}29\,\tfrac{a_c+1}{2}$};
\node[box,anchor=west,text width=4.9cm,font=\scriptsize,minimum height=0.8cm] (d2) at (2.4,-1.1)
     {Softmax độ sắc $\kappa=10$};
\node[box,anchor=west,text width=4.9cm,font=\scriptsize,minimum height=0.8cm] (d3) at (2.4,-2.2)
     {Chiếu Euclid lên đơn hình $\Delta_K$};
\node[box,anchor=west,text width=4.9cm,font=\scriptsize,minimum height=0.8cm] (d4) at (2.4,-3.3)
     {Ánh xạ tuyến tính $[-1,1]\to[0,1]$};
\node[box,anchor=west,text width=4.9cm,font=\scriptsize,minimum height=0.8cm] (d5) at (2.4,-4.725)
     {Pha tham chiếu $\bar\theta_n=-\angle z_n$\\ cộng phần hiệu chỉnh học được};

% ---------- bien vat ly ----------
\node[sbox,anchor=west,fill=blue!10,text width=5.3cm,font=\scriptsize,minimum height=0.8cm] (v1) at (8.1,0)
     {$p_c=\rho_c P_{\max}$, \ $\rho_c\in[0{,}70;\,0{,}99]$};
\node[sbox,anchor=west,fill=blue!10,text width=5.3cm,font=\scriptsize,minimum height=0.8cm] (v2) at (8.1,-1.1)
     {$p_k=(1-\rho_c)P_{\max}w_k$ \ (gần one-hot)};
\node[sbox,anchor=west,fill=blue!10,text width=5.3cm,font=\scriptsize,minimum height=0.8cm] (v3) at (8.1,-2.2)
     {$\eta_k\ge 0$, \ $\textstyle\sum_k \eta_k = 1$};
\node[sbox,anchor=west,fill=blue!10,text width=5.3cm,font=\scriptsize,minimum height=0.8cm] (v4) at (8.1,-3.3)
     {$\beta^T_n\in[0,1]$, \ $\beta^R_n=1-\beta^T_n$};
\node[sbox,anchor=west,fill=blue!10,text width=5.3cm,font=\scriptsize,minimum height=0.8cm] (v5) at (8.1,-4.725)
     {$\theta^T_n,\ \theta^R_n=\bar\theta_n+0{,}25\pi\,a_{\theta}$};

% ---------- noi day: hang ngang, khong cat nhau ----------
\foreach \i/\y in {1/0, 2/-1.1, 3/-2.2, 4/-3.3}{
   \draw[->] (1.2,\y) -- (d\i.west);
}
\draw (1.2,-4.40) -- (1.85,-4.40);
\draw (1.2,-5.05) -- (1.85,-5.05);
\draw (1.85,-4.40) -- (1.85,-5.05);
\draw[->] (1.85,-4.725) -- (d5.west);

\foreach \i in {1,...,5}{ \draw[->] (d\i.east) -- (v\i.west); }

\node[font=\scriptsize,anchor=north west,text width=13.4cm] at (0,-6.05)
  {Mọi ràng buộc công suất, đơn hình và bảo toàn năng lượng được bảo đảm ngay tại bộ giải mã,
   nên không cần phép chiếu hay hạng phạt bổ sung sau khi hành động đã được sinh.};
\end{tikzpicture}
}
\caption{Bộ giải mã hành động có cấu trúc vật lý}
\label{fig:ch3-bo-giai-ma}
\figuresource{Tác giả xây dựng theo tệp \texttt{action.py} (2026)}
\end{figure}

\subsection{Ý nghĩa của việc thu hẹp miền hành động}
\label{ssec:thu-hep-mien}

Bộ giải mã có cấu trúc không bao phủ toàn bộ miền vật lý ở
Chương~\ref{chap:mo-hinh-he-thong}. Đặc biệt, $p_c$ chỉ nằm trong một khoảng con và pha
chỉ nằm trong vùng $\pm\pi/4$ quanh tham chiếu cho từng CSI. Vì vậy, phương pháp đánh đổi
tính tổng quát của miền tìm kiếm lấy khả năng học dễ hơn. Đây là một thiên kiến quy nạp có
chủ đích.

Điều này cũng giải thích tại sao các phương pháp AO có thể đạt tổng tốc độ cao hơn. AO làm
việc trực tiếp trên miền vật lý rộng và có thể thử những nghiệm mà chính sách DRL có cấu
trúc không biểu diễn được. So sánh giữa hai nhóm vì vậy phải được đọc như đánh đổi giữa
miền tìm kiếm và số phép đánh giá, chứ không phải hai bộ giải có cùng không gian tìm kiếm
tuyệt đối.

\section{Phần thưởng, bộ điều khiển QoS và lựa chọn mô hình}
\label{sec:phan-thuong}

\subsection{Chỉ số phần thưởng của môi trường}
\label{ssec:phan-thuong-moi-truong}

Môi trường báo cáo một phần thưởng cơ sở:
\begin{equation}
\label{eq:phan-thuong-moi-truong}
r_{\mathrm{env}} = R_{\mathrm{sum}} - 8V - 8V_2 .
\end{equation}

Trong đó: $R_{\mathrm{sum}}$ là tổng tốc độ; $V$ và $V_2$ là tổng thiếu hụt QoS tuyến tính
và bình phương. Công thức cung cấp một chỉ số duy nhất để quan sát chất lượng hành động.
Khi tất cả UE đạt QoS, $V = V_2 = 0$ nên phần thưởng đúng bằng tổng tốc độ. Khi có vi
phạm, giá trị bị trừ để phản ánh rằng tăng tốc độ bằng cách bỏ rơi UE không phải mục tiêu
mong muốn.

\subsection{Phần thưởng dùng để huấn luyện}
\label{ssec:phan-thuong-huan-luyen}

Trong quá trình huấn luyện, hệ số tuyến tính không cố định mà được điều khiển bởi một biến
đối ngẫu $\lambda_t$:
\begin{equation}
\label{eq:phan-thuong-huan-luyen}
r_t = R_{\mathrm{sum}} - \lambda_t V - 8V_2 .
\end{equation}

Trong đó: $\lambda_t$ là hệ số phạt QoS tuyến tính tại bước huấn luyện $t$; cấu hình khởi
tạo $\lambda_0 = 8$ và giới hạn trong $[4, 64]$. Cách làm này cho phép mức nghiêm khắc với
vi phạm QoS thay đổi theo diễn biến học: nếu chính sách thường xuyên vi phạm, $\lambda_t$
tăng để Critic xem các hành động đó kém hấp dẫn hơn; nếu vi phạm đã thấp, hệ số không cần
tăng vô hạn.

Mức vi phạm được làm trơn bằng trung bình mũ:
\begin{equation}
\label{eq:ema}
\bar{V}_t = \beta_{\mathrm{ema}} \bar{V}_{t-1}
          + (1 - \beta_{\mathrm{ema}}) V_t,
\qquad \beta_{\mathrm{ema}} = 0,95 .
\end{equation}

Trong đó: $\bar{V}_t$ là vi phạm làm trơn và $\beta_{\mathrm{ema}}$ là hệ số ghi nhớ của
trung bình mũ. Phép làm trơn tránh cập nhật hệ số QoS theo một mẫu kênh đơn lẻ có thể quá
tốt hoặc quá xấu; bộ điều khiển nhìn xu hướng vi phạm gần đây nên phản ứng ổn định hơn.

Sau giai đoạn khởi động và theo chu kỳ $1000$ bước, hệ số được cập nhật:
\begin{equation}
\label{eq:cap-nhat-lambda}
\lambda_{t+1} = \mathrm{clip}\!\left[
\lambda_t + \alpha_\lambda(\bar{V}_t - V_{\mathrm{tar}}),\
\lambda_{\min},\ \lambda_{\max} \right],
\end{equation}
trong đó $\alpha_\lambda = 20$ và $V_{\mathrm{tar}} = 0,001$.

Trong đó: $\alpha_\lambda$ là tốc độ cập nhật; $V_{\mathrm{tar}}$ là mức vi phạm mục tiêu;
phép cắt giữ $\lambda$ trong giới hạn cấu hình. Có thể hiểu đây như một bộ điều chỉnh áp
lực QoS: nếu vi phạm trung bình cao hơn mục tiêu, hệ số tăng; nếu thấp hơn, hệ số có thể
giảm. Cơ chế này được áp dụng chung cho ba thuật toán DRL trong giao thức huấn luyện.

\subsection{Lựa chọn điểm lưu mô hình theo tập xác thực}
\label{ssec:chon-diem-luu}

Mô hình không được chọn chỉ bằng phần thưởng hoặc tổng tốc độ. Tại mỗi lần xác thực, hệ
thống tính tổng tốc độ trung bình, tỷ lệ UE đạt QoS, xác suất tất cả UE đạt QoS và vi phạm
trung bình. Một điểm lưu mô hình trước hết phải đáp ứng các ngưỡng QoS đã khai báo. Trong
các điểm lưu khả thi, điểm có tổng tốc độ cao hơn được ưu tiên. Nếu chưa có điểm lưu khả
thi, bộ chọn ưu tiên điểm có khoảng cách chuẩn hóa đến các ngưỡng QoS nhỏ hơn.

Cách chọn này giúp tránh trường hợp mạng đạt tổng tốc độ cao bằng cách hy sinh một số UE.
Đồng thời, việc lựa chọn dựa trên tập xác thực tách biệt giúp giữ tập kiểm thử cho bước
báo cáo cuối thay vì dùng nó để điều chỉnh chính sách.

\section{Kiến trúc mạng nơ-ron}
\label{sec:kien-truc-mang}

\subsection{Mạng Actor xác định của DDPG và TD3}
\label{ssec:actor-xac-dinh}

Mạng Actor của DDPG và TD3 là mạng perceptron đa lớp (Multilayer Perceptron, MLP) kết nối
đầy đủ. Hai lớp ẩn đều có $256$ nút. Sau mỗi lớp kết nối đầy đủ (Fully Connected, FC) ở
phần ẩn là chuẩn hóa theo lớp (Layer Normalization, LayerNorm) và hàm kích hoạt tuyến tính
chỉnh lưu (Rectified Linear Unit, ReLU). Lớp cuối có $d_a$ đầu ra và dùng hàm tang
hyperbolic (Tanh) để đưa mỗi thành phần vào $[-1,1]$.

Lớp FC cuối của Actor được khởi tạo với trọng số và độ lệch rất nhỏ trong khoảng
$[-3 \times 10^{-3},\ 3 \times 10^{-3}]$. Nhờ đó, đầu ra ban đầu gần $0$. Với bộ giải mã
đã xây dựng, hành động gần $0$ tương ứng tỷ lệ công suất chung khoảng $0,845$, chia năng
lượng gần $0,5$, phân bổ công suất riêng ban đầu gần đều và phần hiệu chỉnh pha gần $0$.
Mục đích của cách khởi tạo là bắt đầu từ một vùng có ý nghĩa vật lý thay vì một hành động
cực đoan ngẫu nhiên.

\subsection{Khác biệt giữa Critic của DDPG và TD3}
\label{ssec:critic-ddpg-td3}

DDPG chỉ dùng một Critic $Q(s,a)$. TD3 dùng hai Critic độc lập $Q_1$ và $Q_2$ có cùng cấu
trúc. Hai mạng không chia sẻ trọng số. Khi tạo mục tiêu học, TD3 lấy giá trị nhỏ hơn của
hai Critic. Đây là cơ chế giảm đánh giá quá cao chứ không phải làm mạng sâu hơn.

\subsection{Mạng chính sách và mạng giá trị của PPO}
\label{ssec:mang-ppo}

PPO dùng cùng kiểu MLP hai lớp ẩn $256$ nút với LayerNorm và ReLU, nhưng đầu ra Actor
không đi qua Tanh ngay bên trong MLP. Mạng tạo trung bình $\mu(s)$ của một phân phối
Gaussian. Độ lệch chuẩn được học qua véc-tơ tham số $\log \sigma$. Một mẫu $u$ được lấy từ
Gaussian rồi mới đưa qua Tanh để tạo hành động chuẩn hóa $a = \tanh(u)$. Critic của PPO là
mạng giá trị nhận trạng thái $s$ và xuất $V(s)$, không nhận hành động.

Hình~\ref{fig:ch3-kien-truc-mang} trình bày ba kiến trúc cạnh nhau. Cả ba cột dùng cùng bề
rộng ẩn và cùng thứ tự FC, LayerNorm rồi ReLU, nên khác biệt giữa các thuật toán nằm ở
đầu ra và ở cách tạo hành động chứ không nằm ở năng lực biểu diễn của mạng.

\begin{figure}[htbp]
\centering
\resizebox{\textwidth}{!}{\begin{tikzpicture}[font=\scriptsize,>=Stealth,line width=0.5pt]
\def\bw{3.1}   % be rong mot cot
\def\bh{0.52}  % chieu cao mot khoi

% ---------- cot 1: Actor xac dinh (DDPG / TD3) ----------
\begin{scope}[shift={(0,0)}]
  \node[tit] at (0,0.95) {(a) Actor — DDPG, TD3};
  \node[sm,anchor=north,text width=\bw cm] (a1) at (1.55,0.25)  {Trạng thái $s$ \ ($d_s$)};
  \node[sm,anchor=north,text width=\bw cm] (a2) at (1.55,-0.55) {FC 256 \ $\to$ \ LN + ReLU};
  \node[sm,anchor=north,text width=\bw cm] (a3) at (1.55,-1.35) {FC 256 \ $\to$ \ LN + ReLU};
  \node[sm,anchor=north,text width=\bw cm] (a4) at (1.55,-2.15) {FC $d_a$};
  \node[sm,anchor=north,text width=\bw cm,fill=cChung!10] (a5) at (1.55,-2.95) {Tanh \ $\to$ \ $a\in[-1,1]^{d_a}$};
  \foreach \i/\j in {a1/a2,a2/a3,a3/a4,a4/a5}{ \draw[->] (\i.south) -- (\j.north); }
\end{scope}

% ---------- cot 2: Actor ngau nhien (PPO) ----------
\begin{scope}[shift={(4.9,0)}]
  \node[tit] at (0,0.95) {(b) Actor — PPO};
  \node[sm,anchor=north,text width=\bw cm] (b1) at (1.55,0.25)  {Trạng thái $s$ \ ($d_s$)};
  \node[sm,anchor=north,text width=\bw cm] (b2) at (1.55,-0.55) {FC 256 \ $\to$ \ LN + ReLU};
  \node[sm,anchor=north,text width=\bw cm] (b3) at (1.55,-1.35) {FC 256 \ $\to$ \ LN + ReLU};
  \node[sm,anchor=north,text width=\bw cm] (b4) at (1.55,-2.15) {FC $d_a$ \ $\to$ \ $\mu(s)$};
  \node[sm,anchor=north,text width=\bw cm] (b5) at (1.55,-2.95) {Lấy mẫu $u\sim\mathcal{N}(\mu,\sigma^2)$\\ \tiny $\log\sigma$ là tham số học riêng};
  \node[sm,anchor=north,text width=\bw cm,fill=cChung!10] (b7) at (1.55,-3.75) {Tanh \ $\to$ \ $a=\tanh(u)$};
  \foreach \i/\j in {b1/b2,b2/b3,b3/b4,b4/b5,b5/b7}{ \draw[->] (\i.south) -- (\j.north); }
\end{scope}

% ---------- cot 3: Critic ----------
\begin{scope}[shift={(9.8,0)}]
  \node[tit] at (0,0.95) {(c) Critic};
  \node[sm,anchor=north,text width=\bw cm] (c1) at (1.55,0.25)  {$[s,a]$ \ ($d_s+d_a$) — DDPG/TD3\\ $s$ \ ($d_s$) — PPO};
  \node[sm,anchor=north,text width=\bw cm] (c2) at (1.55,-0.80) {FC 256 \ $\to$ \ LN + ReLU};
  \node[sm,anchor=north,text width=\bw cm] (c3) at (1.55,-1.60) {FC 256 \ $\to$ \ LN + ReLU};
  \node[sm,anchor=north,text width=\bw cm] (c4) at (1.55,-2.40) {FC 1};
  \node[sm,anchor=north,text width=\bw cm,fill=cKenh!10] (c5) at (1.55,-3.20)
       {$Q(s,a)$ — DDPG/TD3\\ $V(s)$ — PPO};
  \foreach \i/\j in {c1/c2,c2/c3,c3/c4,c4/c5}{ \draw[->] (\i.south) -- (\j.north); }
\end{scope}

\node[note,text width=12.8cm] at (0,-4.55)
  {Cả ba cột dùng cùng bề rộng ẩn 256 và cùng thứ tự FC $\to$ LN $\to$ ReLU.
   TD3 dùng hai Critic độc lập có cấu trúc giống cột (c) và lấy giá trị nhỏ hơn khi tạo mục tiêu.
   Đầu ra của cả hai Actor đều đi vào cùng một bộ giải mã hành động vật lý ở Hình 3.1.};
\end{tikzpicture}
}
\caption{Kiến trúc khối của Actor và Critic}
\label{fig:ch3-kien-truc-mang}
\figuresource{Tác giả xây dựng theo tệp \texttt{agents/networks.py} và \texttt{agents/ppo.py} (2026)}
\end{figure}

\section{Nhóm thuật toán học ngoài chính sách}
\label{sec:off-policy}

\subsection{DDPG}
\label{ssec:ddpg}

DDPG dùng Actor xác định $\mu_\phi(s)$, Critic $Q_\psi(s,a)$, mạng đích và bộ nhớ phát
lại. Với một mẫu chuyển trạng thái $(s, a, r, s', d)$, mục tiêu Critic là:
\begin{equation}
\label{eq:muc-tieu-ddpg}
y = r + \gamma (1 - d)\, Q_{\psi^-}\!\big(s',\ \mu_{\phi^-}(s')\big).
\end{equation}

Trong đó: $\phi$ và $\psi$ là tham số Actor và Critic; dấu trừ ở trên chỉ mạng đích;
$\gamma$ là hệ số chiết khấu; $d = 1$ khi bước là cuối giai đoạn. Công thức dùng phần
thưởng hiện tại cộng với giá trị dự kiến của trạng thái tiếp theo. Mạng đích thay đổi chậm
hơn mạng chính, giúp mục tiêu học không chạy quá nhanh theo chính mạng đang được cập nhật.

Critic được học để giảm sai số giữa $Q_\psi(s,a)$ và $y$. Actor được cập nhật theo hướng
làm $Q_\psi(s, \mu_\phi(s))$ tăng. Trong cấu hình cuối, hàm mất mát Critic dùng hàm Huber
và gradient được cắt chuẩn để hạn chế bước cập nhật quá lớn.

\subsection{TD3}
\label{ssec:td3}

TD3 giữ cấu trúc học ngoài chính sách của DDPG nhưng thêm nhiễu làm trơn vào hành động mục
tiêu:
\begin{equation}
\label{eq:nhieu-lam-tron}
\tilde{a}' = \mathrm{clip}\big(\mu_{\phi^-}(s') + \epsilon,\ -1,\ 1\big),
\qquad
\epsilon \sim \mathrm{clip}\big(\mathcal{N}(0, \sigma_{\mathrm{tar}}),\ -c,\ c\big).
\end{equation}

Trong đó: $\tilde{a}'$ là hành động đích đã làm trơn; $\sigma_{\mathrm{tar}}$ là độ lệch
chuẩn nhiễu đích; $c$ là biên cắt nhiễu. Cơ chế này buộc Critic không chỉ đánh giá tốt một
điểm hành động cực hẹp mà còn phải ổn định trong vùng lân cận, có ích khi hàm giá trị xấp
xỉ bằng mạng nơ-ron có thể tạo các đỉnh giả.

Mục tiêu dùng hai Critic:
\begin{equation}
\label{eq:muc-tieu-td3}
y = r + \gamma (1 - d) \min_{i \in \{1,2\}} Q_{\psi_i^-}(s',\ \tilde{a}').
\end{equation}

Trong đó: $Q_{\psi_1^-}$ và $Q_{\psi_2^-}$ là hai Critic đích độc lập; phép lấy nhỏ nhất
chọn giá trị thận trọng hơn. Điều này giảm nguy cơ một Critic đánh giá quá lạc quan rồi
kéo Actor theo một hướng sai; nếu một Critic tạo đỉnh giá trị giả, Critic còn lại có thể
làm mục tiêu thấp xuống.

Actor chỉ cập nhật mỗi hai lần cập nhật Critic theo cấu hình. Hệ số cập nhật mềm của mạng
đích là $\tau = 0,005$:
\begin{equation}
\label{eq:cap-nhat-mem}
\vartheta^- \leftarrow (1 - \tau)\,\vartheta^- + \tau\,\vartheta .
\end{equation}

Trong đó: $\vartheta$ đại diện tham số của Actor hoặc Critic chính và $\vartheta^-$ là
tham số mạng đích. Phép cập nhật làm mạng đích đi theo mạng chính từ từ; với $\tau$ nhỏ,
mục tiêu học thay đổi êm hơn qua các bước, góp phần ổn định Critic.

\subsection{Khám phá và sử dụng lại mẫu}
\label{ssec:kham-pha}

Trong $2.000$ bước khởi động đầu, hành động được lấy ngẫu nhiên đều trong $[-1,1]^{d_a}$
để lấp bộ nhớ phát lại. Sau đó, DDPG và TD3 dùng hành động Actor cộng nhiễu Gaussian khám
phá. Độ lệch chuẩn nhiễu giảm tuyến tính từ $0,12$ xuống $0,03$ trong $100.000$ bước. Với
TD3, nhiễu còn được điều chỉnh theo số chiều hành động bằng hệ số tham chiếu để tránh tổng
năng lượng nhiễu tăng quá mạnh khi $N$ lớn.

Bộ nhớ phát lại có sức chứa $200.000$ mẫu và lô nhỏ kích thước $256$. Mỗi mẫu có thể được
sử dụng nhiều lần ở các bước cập nhật khác nhau. Đây là ưu thế điển hình của nhóm học
ngoài chính sách về hiệu quả sử dụng dữ liệu, nhưng cũng làm Critic phải học từ dữ liệu
được thu dưới các chính sách cũ.

\section{Thuật toán học theo chính sách hiện hành}
\label{sec:on-policy}

PPO thu thập một đoạn tương tác hiện hành rồi cập nhật nhiều vòng trên chính đoạn đó. Tỷ
số xác suất giữa chính sách mới và cũ là:
\begin{equation}
\label{eq:ty-so-ppo}
r_t(\theta) = \frac{\pi_\theta(a_t \mid s_t)}{\pi_{\theta_{\mathrm{old}}}(a_t \mid s_t)} .
\end{equation}

Trong đó: $\pi_\theta$ là chính sách mới; $\pi_{\theta_{\mathrm{old}}}$ là chính sách đã
dùng để thu dữ liệu; $a_t$ và $s_t$ là hành động và trạng thái tại bước $t$. Tỷ số đo mức
chính sách mới thay đổi xác suất của hành động đã quan sát; nếu quá lớn hoặc quá nhỏ, PPO
xem đó là thay đổi quá mạnh so với dữ liệu hiện có.

Hàm mục tiêu cắt của PPO là:
\begin{equation}
\label{eq:muc-tieu-ppo}
L^{\mathrm{clip}}(\theta) = \mathbb{E}_t\Big[
\min\big(r_t(\theta)\hat{A}_t,\
\mathrm{clip}(r_t(\theta),\ 1-\epsilon,\ 1+\epsilon)\hat{A}_t\big)\Big].
\end{equation}

Trong đó: $\hat{A}_t$ là lợi thế ước lượng và $\epsilon$ là biên cắt của PPO. Hàm mục tiêu
giới hạn lợi ích của việc thay đổi chính sách quá xa trong một lần cập nhật; có thể hiểu
đây là cơ chế đi từng bước vừa phải để tránh một lô dữ liệu làm Actor đổi hướng quá mạnh.

Mã nguồn lưu trực tiếp mẫu trước Tanh để tính lại xác suất chính xác hơn trong các vòng
PPO, thay vì khôi phục bằng hàm ngược từ một giá trị đã bị cắt gần $\pm 1$. Đây là chi
tiết triển khai giúp giảm sai lệch số học ở các hành động bão hòa.

\section{Các phương pháp đối chiếu tối ưu truyền thống}
\label{sec:doi-chieu-ao}

\subsection{AnalyticalRIS}
\label{ssec:analytical-ris}

AnalyticalRIS tạo một hành động tham chiếu bằng phân bổ công suất và tốc độ đơn giản, hệ
số năng lượng $0,5$ và căn pha dựa trên tổng kênh ghép tầng. Phương pháp này không lặp tối
ưu tất cả khối nên có chi phí thấp. Vai trò chính là tạo điểm khởi đầu có ý nghĩa cho AO
và cung cấp một đối chiếu về độ trễ rất thấp. Cơ chế của cả ba phương pháp trong mục này
đã được minh họa ở Hình~\ref{fig:ch1-ao-truc-giac}.

\subsection{AO--SCA}
\label{ssec:ao-sca}

AO--SCA dùng AnalyticalRIS làm nghiệm đầu. Khối công suất và khối phân bổ tốc độ chung đều
có cấu trúc đơn hình. Thay vì thay đổi riêng một tọa độ rồi chiếu lại, mã nguồn chuyển
trực tiếp một lượng tài nguyên $\delta$ từ phần tử $i$ sang phần tử $j$:
\begin{equation}
\label{eq:chuyen-khoi}
x'_i = x_i - \delta,
\qquad
x'_j = x_j + \delta .
\end{equation}

Trong đó: $x$ đại diện một véc-tơ trên đơn hình như công suất hoặc $\boldsymbol{\eta}$;
$\delta$ không vượt lượng đang có ở phần tử $i$. Phép chuyển giữ tổng tài nguyên không đổi
ngay trong từng phép thử. Về vật lý, nó giống việc chuyển một ít công suất hoặc tốc độ
luồng chung từ một luồng sang luồng khác, thay vì tăng một biến rồi sửa toàn bộ véc-tơ
bằng phép chiếu.

Bộ giải thử các hướng chuyển khả thi, chọn hướng có độ dốc cải thiện tốt, sau đó tìm nhiều
mức $\delta$ trên hướng đó. Với khối STAR--RIS, đạo hàm theo từng biến được xấp xỉ bằng
sai phân hữu hạn; một bài toán lân cận tạo đề xuất và cơ chế quay lui tăng hệ số lân cận
nếu đề xuất chưa đủ tốt. Sau mỗi vòng, hai đơn hình được đánh bóng lại và kiểm tra khoảng
cách đến điều kiện dừng. Cơ chế này làm AO--SCA là một bộ giải cục bộ có kiểm soát tốt hơn
tại biên đơn hình, nhưng không biến nó thành nghiệm tối ưu toàn cục.

\subsection{AO--Grid}
\label{ssec:ao-grid}

AO--Grid sử dụng tập giá trị rời rạc cho từng khối. Lưới công suất cho phép giá trị bằng
$0$, nhờ đó một luồng riêng có thể được tắt hoàn toàn. Thay vì cập nhật ngay khi gặp cải
thiện đầu tiên, bộ giải xét các ứng viên và chọn bước cải thiện tốt nhất. Đối với biến
STAR--RIS, thuật toán thử cả thứ tự quét tọa độ thuận và nghịch rồi giữ kết quả tốt hơn.
Cách này giảm phụ thuộc vào thứ tự quét nhưng tăng số phép đánh giá.

Điểm quan trọng khi so sánh là AO tối ưu trên miền vật lý rộng hơn miền có cấu trúc của
DRL. Vì vậy, việc AO đạt tổng tốc độ cao hơn không mâu thuẫn với mục tiêu luận văn; nó
phản ánh việc AO dành nhiều phép tính hơn và có quyền truy cập vùng hành động rộng hơn cho
từng CSI.

\section{Ngân hàng kịch bản kênh và giao thức dữ liệu}
\label{sec:ngan-hang-kich-ban}

\subsection{Ngân hàng kịch bản kênh}
\label{ssec:ngan-hang}

Mã nguồn dùng cấu trúc \texttt{ScenarioBank}, trong luận văn gọi là ngân hàng kịch bản
kênh. Với mỗi kích thước $N$, ba ngân hàng độc lập được tạo: $10.000$ kịch bản cho huấn
luyện, $1.000$ cho xác thực và $1.000$ cho kiểm thử. Mỗi ngân hàng lưu đầy đủ
$\mathbf{h}_d$, $\mathbf{g}$, $\mathbf{h}_r$, chỉ báo phía UE và siêu dữ liệu.

Các ngân hàng dùng hạt giống sinh ngẫu nhiên cố định và có mã kiểm tra SHA--256. Một hàm
kiểm tra không trùng lặp duyệt từng kịch bản để bảo đảm ba tập tách biệt. Mục đích không
phải biến kênh tổng hợp thành dữ liệu thật, mà bảo đảm mọi thuật toán được so sánh trên
đúng cùng các mẫu kênh và một lần chạy có thể tái tạo đúng tập dữ liệu.

\subsection{Kích thước và số lần huấn luyện}
\label{ssec:kich-thuoc-huan-luyen}

Năm kích thước STAR--RIS được khảo sát:
\begin{equation}
\label{eq:cac-kich-thuoc-n}
N \in \{16,\ 32,\ 64,\ 96,\ 128\}.
\end{equation}

Trong đó: $N$ là số phần tử của STAR--RIS. Năm mức này tạo năm bậc độ khó từ nhỏ đến lớn.
Khi $N$ tăng, số biến STAR--RIS tăng và đồng thời khả năng tạo độ lợi kênh cũng tăng, vì
vậy có thể đánh giá cả hiệu năng lẫn khả năng mở rộng.

Mỗi thuật toán DRL được huấn luyện với năm hạt giống ngẫu nhiên độc lập trên mỗi $N$, mỗi
lần có ngân sách $100.000$ tương tác môi trường. Tập xác thực được đánh giá định kỳ mỗi
$5.000$ bước trên $1.000$ kịch bản. Tập kiểm thử gồm $1.000$ kịch bản được dùng để báo cáo
chất lượng mô hình đã chọn.

\subsection{Cấu hình chính}
\label{ssec:cau-hinh-chinh}

Bảng~\ref{tab:ch3-cau-hinh} tổng hợp các tham số quan trọng lấy từ cấu hình thực nghiệm.
Các thuật toán DDPG và PPO được cấp các cơ chế ổn định tương ứng như LayerNorm và cắt
gradient để giảm chênh lệch do mức độ hoàn thiện cài đặt.

\begin{table}[htbp]
\centering
\caption{Cấu hình huấn luyện và đánh giá chính}
\label{tab:ch3-cau-hinh}
\footnotesize
\begin{tabularx}{\textwidth}{@{}X l@{}}
\toprule
\textbf{Thành phần} & \textbf{Giá trị} \\
\midrule
Số UE & $K = 4$ \\
Số phần tử STAR--RIS & $N \in \{16, 32, 64, 96, 128\}$ \\
Tổng công suất chuẩn hóa & $P_{\max} = 1$ \\
Công suất tạp âm & $\sigma^2 = 10^{-3}$ \\
Tốc độ QoS tối thiểu & $R_{\min} = 0,5$ bit/s/Hz \\
Độ dài một giai đoạn & $32$ bước \\
Ngân hàng huấn luyện / xác thực / kiểm thử & $10.000 / 1.000 / 1.000$ kịch bản \\
Số hạt giống DRL & $5$ mỗi thuật toán và mỗi $N$ \\
Ngân sách huấn luyện & $100.000$ tương tác mỗi hạt giống và mỗi $N$ \\
Kích thước bộ nhớ phát lại & $200.000$ \\
Kích thước lô nhỏ & $256$ \\
Số nút mỗi lớp ẩn & $256$ \\
Hệ số chiết khấu $\gamma$ & $0,99$ \\
Hệ số cập nhật mềm $\tau$ & $0,005$ \\
Tốc độ học Actor của TD3 và DDPG & $10^{-4}$ \\
Tốc độ học Critic của TD3 và DDPG & $3 \times 10^{-4}$ \\
Tốc độ học PPO & $3 \times 10^{-4}$ \\
Chu kỳ xác thực & $5.000$ bước \\
Nhiễu khám phá của DDPG và TD3 & $0,12 \rightarrow 0,03$ \\
Hàm mất mát Critic của TD3 và DDPG & Huber \\
Cắt chuẩn gradient & $10$ \\
\bottomrule
\end{tabularx}
\tablesource{Tác giả tổng hợp từ cấu hình thực nghiệm (2026)}
\end{table}

\section{Quy trình huấn luyện và đánh giá}
\label{sec:quy-trinh}

Quy trình chung được thể hiện ở Hình~\ref{fig:ch3-quy-trinh}. Tập huấn luyện tạo các tương
tác để cập nhật mạng. Tập xác thực chỉ dùng để chọn điểm lưu mô hình. Sau khi tiêu chí
chọn đã cố định, mô hình được đánh giá trên tập kiểm thử. Các bộ giải AO và AnalyticalRIS
không có giai đoạn huấn luyện; chúng nhận từng CSI kiểm thử và giải trực tiếp.

\begin{figure}[htbp]
\centering
\resizebox{0.95\textwidth}{!}{\begin{tikzpicture}[font=\scriptsize,>=Stealth,line width=0.5pt]
\node[box,text width=2.9cm,minimum height=1.0cm] (b1) at (0,0)
     {Tạo và khóa\\ ba ngân hàng kênh\\ \tiny 10.000 / 1.000 / 1.000};
\node[box,text width=2.9cm,minimum height=1.0cm] (b2) at (4.3,0)
     {Huấn luyện DRL\\ \tiny 100.000 tương tác};
\node[box,text width=2.9cm,minimum height=1.0cm] (b3) at (8.6,0)
     {Xác thực định kỳ\\ \tiny mỗi 5.000 bước};
\node[box,fill=cChung!10,text width=2.9cm,minimum height=1.0cm] (b4) at (8.6,-2.4)
     {Chọn điểm lưu\\ \tiny ưu tiên QoS, sau đó tổng tốc độ};
\node[box,text width=2.9cm,minimum height=1.0cm] (b5) at (4.3,-2.4)
     {Khóa mô hình\\ đã chọn};
\node[box,fill=cKenh!10,text width=2.9cm,minimum height=1.0cm] (b6) at (0,-2.4)
     {Đánh giá kiểm thử\\ \tiny 1.000 kịch bản};
\node[box,text width=2.9cm,minimum height=0.8cm] (b7) at (0,-4.3)
     {Bảng, đồ thị và thống kê};

\draw[->] (b1) -- (b2);
\draw[->] (b2) -- (b3);
\draw[->] (b3) -- (b4);
\draw[->] (b4) -- (b5);
\draw[->] (b5) -- (b6);
\draw[->] (b6) -- (b7);
\draw[->] (b3.west) -- (6.45,0) -- (6.45,-1.25) -- (4.3,-1.25) -- (b2.south);
\node[font=\tiny,anchor=north,fill=white,inner sep=1.5pt] at (5.35,-1.28) {lặp đến hết ngân sách};

\node[box,text width=3.4cm,minimum height=0.8cm] (ao) at (5.2,-4.3)
     {AO--SCA, AO--Grid, AnalyticalRIS\\ \tiny giải trực tiếp từng kịch bản};
\draw[->] (ao.west) -- (b7.east);
\draw[->] (b1.south) -- (0,-0.9) -- (-1.95,-0.9) -- (-1.95,-4.3) -- (b7.west);
\node[font=\tiny,anchor=east,rotate=90] at (-2.02,-2.4) {ngân hàng kiểm thử};
\end{tikzpicture}
}
\caption{Quy trình huấn luyện, xác thực và kiểm thử}
\label{fig:ch3-quy-trinh}
\figuresource{Tác giả xây dựng theo giao thức mã nguồn (2026)}
\end{figure}

\begin{algorithm}[htbp]
\caption{Quy trình huấn luyện chung của nhóm DRL}
\label{alg:huan-luyen}
\begin{algorithmic}[1]
\Require Ngân hàng kênh huấn luyện $\mathcal{B}_{\mathrm{tr}}$, ngân hàng xác thực
         $\mathcal{B}_{\mathrm{val}}$, cấu hình, hạt giống ngẫu nhiên
\Ensure  Điểm lưu mô hình được lựa chọn theo tập xác thực
\State Khởi tạo mạng Actor hoặc mạng chính sách, Critic hoặc mạng giá trị và các bộ nhớ cần thiết
\State Đánh giá mô hình khởi tạo trên $\mathcal{B}_{\mathrm{val}}$ nếu cấu hình yêu cầu
\For{$t \gets 1$ \textbf{to} $100000$}
  \State Lấy một kịch bản CSI từ $\mathcal{B}_{\mathrm{tr}}$ và tạo trạng thái chuẩn hóa
  \State Sinh hành động chuẩn hóa theo DDPG, TD3 hoặc PPO
  \State Giải mã thành $\mathbf{p}, \boldsymbol{\eta}, \boldsymbol{\beta},
         \boldsymbol{\theta}^T, \boldsymbol{\theta}^R$
  \State Tính kênh hiệu dụng, tốc độ, vi phạm QoS và phần thưởng huấn luyện
  \State Cập nhật bộ điều khiển QoS và tham số thuật toán theo cơ chế tương ứng
  \If{$t$ đến chu kỳ xác thực}
    \State Đánh giá chính sách xác định trên $\mathcal{B}_{\mathrm{val}}$
    \State Lưu điểm ứng viên và cập nhật điểm tốt nhất theo tiêu chí QoS trước, tổng tốc độ sau
  \EndIf
\EndFor
\end{algorithmic}
\end{algorithm}

\section{Khả năng tái lập}
\label{sec:kha-nang-tai-lap}

Để một kết quả có thể kiểm tra, mỗi lần chạy cần đi kèm cấu hình, hạt giống ngẫu nhiên, mã
kiểm tra ngân hàng kênh, phiên bản mã nguồn, điểm lưu mô hình, dữ liệu huấn luyện và xác
thực, bảng kiểm thử cùng thông tin phần cứng và phần mềm. Mã nguồn lưu các thông tin này
trong các tệp kê khai, tệp đo thời gian và bảng dữ liệu. Các ngân hàng kênh được kiểm tra
số lượng, kích thước, hạt giống, tính không trùng và mã kiểm tra trước khi chạy.

Luận văn phân biệt rõ ba nguồn bất định. Thứ nhất là biến thiên kênh giữa các kịch bản.
Thứ hai là biến thiên huấn luyện giữa các hạt giống của DRL. Thứ ba là biến thiên thời
gian chạy khi đo độ trễ. Việc tách các nguồn này giúp tránh diễn giải sai một khoảng tin
cậy theo kịch bản như thể đó là độ bất định của toàn bộ thuật toán học.

\section{Kết luận chương}
\label{sec:ch3-ket-luan}

Chương~\ref{chap:phuong-phap} đã trình bày đầy đủ từ trạng thái CSI đến biến vật lý cuối
cùng. Điểm trung tâm là bộ giải mã hành động có cấu trúc, trong đó công suất chung được
giới hạn trong vùng ưu tiên, công suất riêng dùng softmax độ sắc cao, tốc độ luồng chung
được chiếu lên đơn hình, phân chia năng lượng được bảo đảm bằng ánh xạ tuyến tính và pha
được học như phần hiệu chỉnh quanh tham chiếu kênh. DDPG, TD3 và PPO dùng chung bộ giải mã
này nhưng khác cơ chế học. Chương~\ref{chap:ket-qua} sử dụng cùng giao thức dữ liệu để so
sánh nhóm DRL với các phương pháp AO về chất lượng nghiệm, QoS và độ trễ.
